\chapter{一笔 strict replay 订单的完整生命周期}

\section{本章要解决的困惑}

fast 输出 entries/exits。strict replay 输出 event log。要理解 Runner，最好的方式不是读全文件，而是跟踪一笔订单。

一笔订单的链条：

\begin{lstlisting}
market event
-> Bot feature update
-> order_intent
-> order_scheduled
-> order_arrived
-> order_ack / order_reject
-> fill_created
-> portfolio_state_on_change
\end{lstlisting}

\section{market event}

Runner 从 canonical market truth 中读 quote、trade、L2，然后发给 Bot。Bot 只能看到已经发送给它的 public stream。

\begin{lstlisting}
market_quote
market_trade
market_l2_update
market_decision_frame
\end{lstlisting}

在 panel sparse fast clock 模式下，Runner 可以发送 \code{market_decision_frame}，但 market truth 仍来自 canonical 数据。

\section{Bot feature update}

Bot 更新：

\begin{lstlisting}
rolling trades -> TFI
quote -> spread and frames
closed outcomes -> R5
shadow policy -> signal
\end{lstlisting}

如果触发 entry，就产生 order intent。

\section{order\_intent}

order intent 是策略意图，不是成交事实。它应该记录：

\begin{lstlisting}
side
qty or exposure intent
observed_seq
observed_local_ts_us
feature_snapshot
\end{lstlisting}

这个 intent 绑定到 Bot 当时看到的 market event。

\section{order\_scheduled 和 order\_arrived}

Runner 根据 latency 计算到达时间：

\[
target\_arrival\_local\_ts
=observed\_local\_ts+latency\_us.
\]

timer arrival 模式下，Runner 用这个虚拟目标时间选 arrival quote。next-event 模式则等第一个更晚市场事件，更保守。

\section{ack/reject 和 fill}

如果订单合法，Runner ack；否则 reject。当前 taker IOC 主要关心到达时立即成交：

\begin{itemize}[leftmargin=2em]
  \item buy fill price = arrival ask；
  \item sell fill price = arrival bid。
\end{itemize}

fill 之后，portfolio 更新 position、cash、equity、leverage。

\section{event log 怎么审}

\code{events.ndjson} 一行一个事件。每行有 envelope：

\begin{lstlisting}
event_id
run_id
event_type
replay_seq
replay_ts
wall_ts_ms
source
payload
\end{lstlisting}

审计一笔订单时，从 \code{fill_created} 往前追：

\begin{lstlisting}
fill_created
<- order_ack
<- order_arrived
<- order_scheduled
<- order_intent
<- strategy_feature_checkpoint or market event
\end{lstlisting}

\section{Monitor 的角色}

Monitor 只把这些事件按 intent/order 聚合成可视化因果链。它不能下单，不能撤单，不能创造 missing fact。

\warningline{Monitor 只展示，不创造事实。}

\section{一笔 long order 的完整数字例子}

假设 Bot 在 observed 时刻看到：

\begin{lstlisting}
observed_seq = 1000
observed_bid = 0.1120
observed_ask = 0.1124
side = buy
target_exposure = 0.50
\end{lstlisting}

策略意图是：

\begin{lstlisting}
order_intent:
  intent_id = I-1000
  side = buy
  observed_seq = 1000
  observed_quote = 0.1120 / 0.1124
\end{lstlisting}

Runner 加入 latency queue：

\[
arrival\_ts=observed\_ts+latency\_us.
\]

到达时盘口变成：

\begin{lstlisting}
arrival_bid = 0.1121
arrival_ask = 0.1127
\end{lstlisting}

因为 side 是 buy，fill price 是 arrival ask：

\[
fill\_price=0.1127.
\]

这时 event chain 应该能看到：

\begin{lstlisting}
order_intent(observed ask 0.1124)
order_scheduled(target arrival ts)
order_arrived(arrival ask 0.1127)
order_ack
fill_created(fill price 0.1127)
portfolio_state_on_change
\end{lstlisting}

如果 Monitor 里显示 observed quote 是 0.1124 ask，但 fill price 是 0.1127，这不一定是 bug；它可能是 latency 后 arrival quote 变差。

\section{一笔 short order 的完整数字例子}

short 正好相反。observed：

\begin{lstlisting}
observed_bid = 0.1120
observed_ask = 0.1124
side = sell
\end{lstlisting}

arrival：

\begin{lstlisting}
arrival_bid = 0.1118
arrival_ask = 0.1122
\end{lstlisting}

sell taker 打 bid：

\[
fill\_price=0.1118.
\]

如果 60 秒后要平 short，买回会打 exit ask。strict replay 会把开仓 fill 和后续 portfolio 都记录下来；fast fixed60 则通常只从 quote index 直接算读数。

\section{为什么 NDJSON 合适}

\code{events.ndjson} 一行一个 JSON event。这样做有三个好处：

\begin{enumerate}[leftmargin=2em]
  \item 可以边 replay 边写，不需要整段跑完才输出；
  \item 大文件可以逐行 streaming parse；
  \item 每个事件都有 envelope，方便按 \code{intent_id} 或 \code{order_id} 审计。
\end{enumerate}

一行简化 event 像这样：

\begin{lstlisting}
{
  "event_type": "fill_created",
  "replay_seq": 1010,
  "replay_ts": "...",
  "payload": {
    "intent_id": "I-1000",
    "side": "buy",
    "fill_price": 0.1127,
    "arrival_bid": 0.1121,
    "arrival_ask": 0.1127
  }
}
\end{lstlisting}

这不是给策略读的 label，而是给人和 Monitor 审计的事实记录。

\section{strict replay 的常见误读}

\begin{longtable}{p{0.30\textwidth}p{0.56\textwidth}}
\toprule
误读 & 正确理解 \\
\midrule
fill price 应该等于 observed ask & taker fill 用 arrival quote，不是 observed quote \\
Monitor 能修正成交 & Monitor 只展示 event log，不能创造事实 \\
Runner 可以看 fast label & Runner 不读 label，只用市场和订单 \\
ack 等于赚钱 & ack 只表示订单合法，不表示策略收益为正 \\
portfolio state 是策略信号 & portfolio 是成交后账户事实，不是 alpha 输入 \\
\bottomrule
\end{longtable}

\section{手动检查点}

拿一笔 strict run，检查：

\begin{enumerate}[leftmargin=2em]
  \item observed quote 和 arrival quote 是否记录；
  \item fill price 是否符合 side；
  \item arrival quote lag 是否为 0 或有解释；
  \item portfolio state 是否在 fill 后变化；
  \item event log 是否能从 fill 追到 intent。
\end{enumerate}

\checkpoint{strict replay 的价值不是更会赚钱，而是证明订单在本地交易所里是可审计的。}
